\documentclass[12pt,oneside,a4paper]{book}

%==================================================================
% METADATOS
%==================================================================
\newcommand{\universidad}{Universidad de Buenos Aires}
\newcommand{\facultad}{Facultad de Derecho}
\newcommand{\materia}{Derecho Constitucional Profundizado}
\newcommand{\titulo}{Control de Constitucionalidad, Fallos Fundacionales y Sistema Difuso}
\newcommand{\autor}{Isaac Edgar Camacho Ocampo}
\newcommand{\anio}{Año académico 2024--2025}

%==================================================================
% PAQUETES
%==================================================================
\usepackage[utf8]{inputenc}
\usepackage[T1]{fontenc}
\usepackage[provide=*,spanish]{babel}

\usepackage[
    a4paper,
    top=2.5cm,
    bottom=2.5cm,
    left=3cm,
    right=2.5cm,
    headheight=15pt
]{geometry}

\usepackage{amsmath}
\usepackage{amssymb}
\usepackage{mathtools}

\usepackage{graphicx}
\usepackage{float}
\usepackage{caption}
\usepackage{subcaption}

\usepackage{booktabs}
\usepackage{array}
\usepackage{longtable}
\usepackage{tabularx}

\usepackage{enumitem}

\usepackage{xcolor}
\usepackage[most]{tcolorbox}

\usepackage{fancyhdr}
\usepackage{ragged2e}

\usepackage[
    colorlinks=true,
    linkcolor=blue!50!black,
    citecolor=green!40!black,
    urlcolor=blue!60!black,
    pdftitle={Apuntes de Derecho Constitucional Profundizado},
    pdfauthor={Isaac Edgar Camacho Ocampo},
    pdfsubject={Apuntes universitarios},
    bookmarks=true
]{hyperref}

%==================================================================
% CONFIGURACIÓN
%==================================================================
\graphicspath{
    {./img/}
    {./graficos/}
    {./figuras/}
}

\setcounter{tocdepth}{2}

\setlength{\parindent}{0pt}
\setlength{\parskip}{0.5em}

\setlist[itemize]{
    topsep=0.3em,
    itemsep=0.2em
}

\setlist[enumerate]{
    topsep=0.3em,
    itemsep=0.2em
}

\clubpenalty=10000
\widowpenalty=10000

%==================================================================
% COMANDOS
%==================================================================
\newcommand{\abs}[1]{\left|#1\right|}
\newcommand{\norm}[1]{\left\lVert#1\right\rVert}
\newcommand{\vect}[1]{\mathbf{#1}}
\newcommand{\deriv}[2]{\frac{d#1}{d#2}}
\newcommand{\pderiv}[2]{\frac{\partial#1}{\partial#2}}

% Separador de filas del cuadro Cornell
\newcommand{\cornellsep}{\\ \addlinespace \cmidrule{1-2} \addlinespace}

%==================================================================
% ENTORNOS / CAJAS
%==================================================================
\newtcolorbox{definition}[1]{
    enhanced,
    breakable,
    title={Definición: #1},
    fonttitle=\bfseries,
    colback=blue!3,
    colframe=blue!50!black,
    boxrule=0.8pt,
    arc=2mm
}

\newtcolorbox{important}{
    enhanced,
    breakable,
    title={Importante},
    fonttitle=\bfseries,
    colback=yellow!8,
    colframe=orange!70!black,
    boxrule=0.8pt,
    arc=2mm
}

\newtcolorbox{example}[1]{
    enhanced,
    breakable,
    title={Ejemplo: #1},
    fonttitle=\bfseries,
    colback=green!4,
    colframe=green!40!black,
    boxrule=0.8pt,
    arc=2mm
}

\newtcolorbox{formula}[1]{
    enhanced,
    breakable,
    title={Fórmula / Regla: #1},
    fonttitle=\bfseries,
    colback=purple!4,
    colframe=purple!50!black,
    boxrule=0.8pt,
    arc=2mm
}

\newtcolorbox{exercise}[1]{
    enhanced,
    breakable,
    title={Ejercicio: #1},
    fonttitle=\bfseries,
    colback=gray!5,
    colframe=gray!60!black,
    boxrule=0.8pt,
    arc=2mm
}

\newtcolorbox{note}{
    enhanced,
    breakable,
    title={Nota},
    fonttitle=\bfseries,
    colback=white,
    colframe=gray!50,
    boxrule=0.6pt,
    arc=2mm
}

%==================================================================
% CONFIGURACIÓN FINAL
%==================================================================
\pagestyle{fancy}
\fancyhf{}

\fancyhead[L]{\nouppercase{\leftmark}}
\fancyhead[R]{\materia}

\fancyfoot[C]{\thepage}

\renewcommand{\headrulewidth}{0.4pt}
\renewcommand{\footrulewidth}{0pt}

\fancypagestyle{plain}{
    \fancyhf{}
    \fancyfoot[C]{\thepage}
    \renewcommand{\headrulewidth}{0pt}
}

%==================================================================
\begin{document}
%==================================================================

%------------------------------------------------------------------
% PORTADA
%------------------------------------------------------------------
\begin{titlepage}

\thispagestyle{empty}

\begin{center}

\vspace*{1cm}

{\large \textsc{\universidad}}

\vspace{0.2cm}

{\large \textsc{\facultad}}

\vspace{3cm}

{\Huge \textbf{\materia}}

\vspace{1cm}

{\LARGE \titulo}

\vspace{0.8cm}

\textit{Comprender es el primer paso para convertir el estudio en conocimiento.}

\vspace{1.2cm}

\rule{0.70\textwidth}{0.5pt}

\vspace{0.5cm}

{\large Apuntes de estudio}

\vspace{0.5cm}

\rule{0.70\textwidth}{0.5pt}

\vfill

{\large \autor}

\vspace{0.3cm}

{\large \anio}

\end{center}

\vfill

\begin{flushleft}
\textit{Este es un modesto aporte a los estudiantes de ``\materia'' no pretende sustituir el material oficial de la materia.}
\end{flushleft}

\end{titlepage}

%------------------------------------------------------------------
% ÍNDICE
%------------------------------------------------------------------
\tableofcontents

%==================================================================
% CAPÍTULO 1
%==================================================================
\chapter{Fallo Marbury v. Madison (1803)}
\markboth{Fallo Marbury v. Madison (1803)}{}

El estudio de los fallos fundacionales del control de constitucionalidad parte de una situación hipotética: un ciudadano es detenido injustamente por publicar una caricatura crítica en redes sociales. Surgen entonces preguntas como ante qué juez acudir y si es posible ir directamente a la Corte Suprema. Esas cuestiones son las que se plantearon en 1887 en el caso Sojo y en 1803 en el caso Marbury. El hilo conductor de todo el material es uno solo: \textbf{quién tiene la última palabra sobre la Constitución y cómo se ejerce ese poder}.

\begin{note}
Los apuntes señalan la relevancia profesional de estos contenidos:
\begin{itemize}
    \item Cualquier norma o acto estatal puede ser cuestionado ante los tribunales si viola la Constitución; conocer el sistema permite identificar cuándo hay una vulneración de derechos y cómo atacarla judicialmente.
    \item El recurso extraordinario federal es la herramienta que permite llegar a la Corte Suprema.
    \item Los conceptos de jurisdicción y competencia federal son decisivos para determinar ante qué tribunal demandar o defenderse.
    \item La división de poderes es el límite que impide que el Ejecutivo o el Legislativo se extralimiten.
    \item Los fallos Marbury v. Madison y Sojo son citados constantemente en jurisprudencia y doctrina.
\end{itemize}
\end{note}

\section{Contexto político: Adams y Jefferson}

\subsection{La legitimidad de los nombramientos}

El primer punto de análisis es si un presidente saliente tiene legitimidad para nombrar jueces. El presidente Adams, a punto de dejar el poder, había nombrado a \textbf{42 jueces} que todavía no habían asumido. Al haber perdido la elección, el mandato que le restaba era un mandato \textbf{remanente} y la opinión del pueblo era la opuesta. La pregunta que se plantea es si quien acaba de perder las elecciones tiene la legitimidad necesaria para nombrar jueces, aunque, como se señaló en el debate, el mandato vigente lo autorice formalmente.

Cuando Jefferson ganó las elecciones, quedaban esos nombramientos pendientes. Jefferson sostenía que no los nombraría porque había ganado la elección, mientras que se le respondía que el acto de nombramiento era válido.

\subsection{La validez del acto de nombramiento}

El acto de nombramiento \textbf{existió} en la forma en que el término se utilizaba: era un \textbf{acto de Estado}. Lo único que le faltaba era la rubricación, pero existió. La posición de Jefferson (``yo gané en las elecciones'') no figura expresamente en el fallo, sino que forma parte del contexto.

De este modo, Jefferson se creía con derecho no solo a no nombrar a los jueces, sino también a \textbf{interpretar la Constitución}, es decir, a ejercer esa función. Ese era el dilema central del fallo.

\subsection{La ambigüedad de la Constitución norteamericana}

La Constitución de los Estados Unidos no era clara, y aún hoy no lo es, al establecer cuáles son las funciones propias de la Corte Suprema de Justicia y cuáles las de los demás tribunales ordinarios inferiores.

En el caso se presentaban dos cuestiones a la vez: la \textbf{cuestión política} de un gobierno que salía y otro que entraba, y la cuestión de \textbf{jurisdicción y competencia}. En los Estados Unidos existen permanentemente pujas de poder entre lo local y lo federal.

\section{El federalismo norteamericano en comparación con el argentino}

Para un norteamericano promedio, la Constitución nacional resulta menos importante que las leyes locales, porque son estas las que rigen la vida cotidiana de cada persona. Las provincias o los estados tienen mucho poder. Se rigen por códigos propios y no existen códigos de fondo nacionales como los que hay en Argentina.

\begin{example}{Diferencias entre estados en los Estados Unidos}
Existe un código civil y comercial de Luisiana y otro de Utah, y también una ley penal de Utah y otra de Luisiana. No hay un código penal nacional. Por eso en algunos estados existe la pena de muerte y en otros no, y en algunos el aborto está penalizado y en otros no.
\end{example}

\begin{table}[H]
\centering
\begin{tabularx}{\textwidth}{
    >{\raggedright\arraybackslash}p{0.22\textwidth}
    >{\raggedright\arraybackslash}X
    >{\raggedright\arraybackslash}X
}
\toprule
\textbf{Materia} & \textbf{Estados Unidos} & \textbf{Argentina} \\
\midrule
Derecho civil y comercial & Cada estado tiene su propio código (por ejemplo, Luisiana y Utah). & Un mismo Código Civil y Comercial se aplica en todo el país. \\
\addlinespace
Derecho penal & Cada estado tiene su propia ley penal; no hay un código penal nacional. & Un mismo Código Penal se aplica en todo el país. \\
\bottomrule
\end{tabularx}
\caption{Códigos de fondo en Estados Unidos y en Argentina.}
\end{table}

Argentina también presenta diferencias entre provincias: por ejemplo, cada provincia tiene su propio \textbf{código de procedimiento}. Sin embargo, siempre existen leyes federales, como el Código Civil y Comercial y el Código Penal, que se aplican igual en todo el país.

\section{La vía procesal elegida por Marbury}

Marbury era uno de los 42 jueces nombrados que reclamaba ser puesto en funciones. Su planteo no se hizo ante un juez de primera instancia del Distrito de Columbia, donde está Washington, sino \textbf{directamente ante la Corte}, para reclamar que se los nombrara en el cargo.

El dilema era si habían hecho bien en ir directamente a la Corte o si tendrían que haber acudido a un juez ordinario, es decir, si correspondía reclamar el nombramiento ante el Poder Judicial o en otro lugar. Nombrar jueces, ya sean de primera instancia, de segunda instancia o de la Corte, es un \textbf{acto federal complejo}, porque intervienen dos poderes: el Poder Ejecutivo, que nombra, y el Poder Legislativo, a través del Senado, que presta el aval para esos nombramientos.

\section{La figura de John Marshall}

Marshall era el presidente de la Corte y, irónicamente, había sido Secretario de Estado. Se lo describe como más perspicaz y astuto, y se señala que, visto retrospectivamente, actuó como correspondía: actuó como la Corte.

En lugar de elegir el camino fácil (``sí, hay que nombrarlos porque el acto existió''), Marshall observó que la Constitución norteamericana no es clara respecto de cuándo se llega directamente a la Corte y cuándo se debe pasar por los tribunales inferiores. Examinó entonces el texto constitucional para determinar qué estaba en conflicto.

\subsection{El conflicto normativo}

El conflicto normativo central del fallo se daba entre:
\begin{itemize}
    \item la \textbf{Judiciary Act} (ley judicial número 13), que establecía que había que nombrar a estos jueces; y
    \item el \textbf{artículo 3 de la Constitución norteamericana}, que no establecía claramente qué casos van directamente a la Corte y qué casos van por apelación o por la vía ordinaria, es decir, la de primera instancia.
\end{itemize}

\section{La lógica de Marshall: las tres preguntas}

Para resolver el caso, Marshall desarrolló un razonamiento conocido como la \textbf{``lógica de Marshall''}, que se estructura en tres preguntas. Para profundizar en este razonamiento se remite a los manuales de Sola y de Sabsay.

% TODO: contenido incompleto o ambiguo en las notas originales (la segunda pregunta se formula de manera algo distinta en el cuadro de la clase y en el resumen Cornell del apunte original)
\begin{table}[H]
\centering
\begin{tabularx}{\textwidth}{
    >{\centering\arraybackslash}p{0.09\textwidth}
    >{\raggedright\arraybackslash}X
    >{\raggedright\arraybackslash}X
}
\toprule
\textbf{Pregunta} & \textbf{Planteo} & \textbf{Respuesta de Marshall} \\
\midrule
\textbf{1.ª} & \textit{¿Existió el acto administrativo de nombramiento?} & \textbf{Sí.} El acto existió, aunque faltaba la rubricación. En principio, asiste derecho a Marbury para reclamar el puesto. \\
\addlinespace
\textbf{2.ª} & \textit{¿Tiene viabilidad el reclamo? ¿Vinieron al lugar adecuado?} & \textbf{Sí.} El Poder Judicial es el lugar correcto para que se interprete si se violó o no el derecho al nombramiento. \\
\addlinespace
\textbf{3.ª} & \textit{¿Puede el Poder Judicial emitir un mandamus para colocarlos en el cargo?} & \textbf{No.} Hacerlo implicaría interferir con la división de poderes e invadir funciones propias del Poder Ejecutivo. Sin embargo, se logra consolidar al Poder Judicial como único intérprete de la Constitución. \\
\bottomrule
\end{tabularx}
\caption{La lógica de Marshall.}
\end{table}

\subsection{La respuesta definitiva: el Poder Judicial no nombra, pero interpreta}

Ante la tercera pregunta, la respuesta fue negativa: nombrar excedería las funciones del Poder Judicial, porque el acto de nombrar empleados y funcionarios públicos no es propio del Poder Judicial sino del \textbf{Poder Ejecutivo}, y la Corte estaría invadiendo su ámbito y usurpando sus funciones. Por lo tanto, si bien a los reclamantes les asistía el derecho, nombrarlos excedía las funciones de la Corte.

Sin embargo, al no poder nombrarlos, el Poder Judicial \textbf{se autorrestringió y, paradójicamente, ganó más}. Le puso límites al presidente norteamericano y se afirmó frente a él. La duda de Jefferson, quien consideraba que él también podía interpretar la Constitución porque lo había elegido el pueblo, quedó resuelta: a partir de ese momento, \textbf{el único órgano capacitado para interpretar y decir lo que es la ley es el Poder Judicial}.

\subsection{``Un gobierno de leyes y no de hombres''}

Todas las personas están por debajo de la ley y nadie está por encima de ella. Con esta idea, la Corte le decía a Jefferson que él también estaba por debajo de la ley, aunque fuera el presidente. La fuerza de ese mensaje se comprende por el contexto: Jefferson atacaba a los miembros de la Corte, a quienes denominaba en inglés \textit{``the nine oldies''} (``los nueve viejitos''), de manera despectiva, como si dijera: ``¿estos nueve viejitos me van a poner a mí un límite?''. Ese era el trasfondo del conflicto.

\begin{important}
\textbf{Frase célebre del fallo:} ``El gobierno de los Estados Unidos es un gobierno de leyes y no de hombres'' (hoy sería: un gobierno de leyes y no de personas). En consecuencia, también el presidente debe acatar la ley, aunque se crea representante del pueblo.
\end{important}

\section{``No tenemos ni la bolsa ni la espada'': el poder más frágil}

No es cierto que el Poder Judicial sea el más fuerte de los tres poderes. Según el fallo, es el \textbf{más débil}, porque no tiene ni la bolsa ni la espada:
\begin{itemize}
    \item la \textbf{bolsa} representa los fondos del Estado, es decir, el ingreso y la recaudación. Quien los controla es, en definitiva, el Poder Legislativo, que establece los impuestos y las tasas;
    \item la \textbf{espada} representa las fuerzas de seguridad, es decir, el monopolio del uso de la fuerza. El Estado ejerce la violencia legal porque tiene ese monopolio.
\end{itemize}

\begin{table}[H]
\centering
\begin{tabularx}{\textwidth}{
    >{\raggedright\arraybackslash}p{0.19\textwidth}
    >{\raggedright\arraybackslash}X
    >{\raggedright\arraybackslash}X
}
\toprule
\textbf{Poder} & \textbf{Recurso / fortaleza} & \textbf{Debilidad} \\
\midrule
Poder Ejecutivo & Monopolio del uso de la fuerza (la ``espada''). & Sometido a la ley y al control judicial. \\
\addlinespace
Poder Legislativo & Control de los fondos públicos (la ``bolsa''). & Sus leyes pueden ser declaradas inconstitucionales. \\
\addlinespace
Poder Judicial & Control de constitucionalidad y última palabra. & No tiene ni la bolsa ni la espada: es el poder más frágil. \\
\bottomrule
\end{tabularx}
\caption{Los tres poderes según Marshall.}
\end{table}

\section{Importancia del fallo}

A partir de este fallo, el Poder Judicial comienza a ejercer el \textbf{control de constitucionalidad}, y el Ejecutivo no lo ejerce: se encuentra por debajo de la ley, lo que expresa la idea de un gobierno de leyes y no de hombres. Esa es la \textbf{autonomía propia} que el Poder Judicial logró con este fallo en los Estados Unidos, y que Argentina \textbf{copió tal cual}. Ahí radica la importancia del fallo Marbury.

\section{Cuadro Cornell: Marbury v. Madison}

\begin{longtable}{@{}>{\RaggedRight\arraybackslash}p{0.27\textwidth}>{\RaggedRight\arraybackslash}p{0.64\textwidth}@{}}
\toprule
\textbf{Preguntas / palabras clave} & \textbf{Notas / respuestas} \\
\midrule
\endfirsthead
\toprule
\textbf{Preguntas / palabras clave} & \textbf{Notas / respuestas} \\
\midrule
\endhead
\bottomrule
\endfoot
\bottomrule
\endlastfoot

\textbf{¿Qué es Marbury v. Madison y por qué es fundacional?} &
Es el fallo de la Corte Suprema de los Estados Unidos de 1803 en el que se ejerce el control de constitucionalidad judicial. John Marshall, presidente de la Corte, determinó que el Poder Judicial es el único intérprete de la Constitución y que, en un gobierno de leyes, ningún presidente, incluido Jefferson, está por encima de la ley. Argentina copió tal cual la autonomía que el Poder Judicial logró con este fallo.
\cornellsep

\textbf{¿Cuál era el contexto político del caso?} &
El presidente saliente Adams había nombrado 42 jueces que todavía no habían asumido, con un mandato remanente y contra la opinión del pueblo. El acto existió como acto de Estado; solo faltaba la rubricación. Jefferson, quien había ganado las elecciones, se negó a nombrarlos y se creía con derecho a interpretar la Constitución. Marbury, uno de los 42, reclamó directamente ante la Corte.
\cornellsep

\textbf{¿Qué normas estaban en conflicto?} &
La Judiciary Act (ley judicial número 13), que establecía que había que nombrar a estos jueces, y el artículo 3 de la Constitución norteamericana, que no establecía claramente qué casos van directamente a la Corte y cuáles por apelación o por la vía ordinaria de primera instancia.
\cornellsep

\textbf{¿Qué es la ``lógica de Marshall''?} &
Es el razonamiento de tres preguntas con el que Marshall resolvió el caso: 1) ¿existió el acto de nombramiento? (sí); 2) ¿tiene viabilidad el reclamo, es decir, se acudió al lugar adecuado? (sí, el Poder Judicial es el lugar correcto para interpretar si se violó el derecho al nombramiento); 3) ¿puede el Poder Judicial emitir un mandamus para colocarlos en el cargo? (no, interferiría con la división de poderes e invadiría funciones del Ejecutivo).
\cornellsep

\textbf{¿Por qué el Poder Judicial, al negarse a nombrar, ganó más poder?} &
Porque, al autorrestringirse, le puso límites al presidente y resolvió la duda de Jefferson sobre si él también podía interpretar la Constitución: a partir de entonces el único órgano capacitado para interpretar y decir lo que es la ley es el Poder Judicial.
\cornellsep

\textbf{¿Qué significa ``un gobierno de leyes y no de hombres''?} &
Que todas las personas están por debajo de la ley y nadie está por encima de ella. La Corte le decía a Jefferson que también él debía acatar la ley, aunque se creyera representante del pueblo y fuera presidente.
\cornellsep

\textbf{``No tenemos ni la bolsa ni la espada'': ¿qué significa?} &
El Poder Judicial es el más débil de los tres poderes: no controla los fondos del Estado (la bolsa, en manos del Legislativo, que establece impuestos y tasas) ni el monopolio del uso de la fuerza (la espada, propia del Estado). Su fuerza es el control de constitucionalidad y la última palabra sobre la ley.
\cornellsep

\textbf{¿En qué se diferencian el federalismo norteamericano y el argentino en materia de códigos?} &
En Estados Unidos cada estado tiene sus propios códigos civil, comercial y penal (por ejemplo, Luisiana y Utah); por eso, entre otras cosas, la pena de muerte existe en algunos estados y no en otros. En Argentina los códigos de fondo son nacionales (Código Civil y Comercial y Código Penal), aunque cada provincia tiene su código de procedimiento. \\
\midrule

\multicolumn{2}{@{}p{\dimexpr0.91\textwidth+2\tabcolsep\relax}@{}}{
\textbf{Resumen:} Marbury v. Madison enfrentó a un presidente saliente que nombró 42 jueces con el nuevo presidente que se negó a nombrarlos. Marshall, mediante sus tres preguntas, reconoció el derecho de Marbury pero sostuvo que el Poder Judicial no podía nombrarlo sin invadir funciones del Ejecutivo. Al autorrestringirse, el Poder Judicial se consolidó como único intérprete de la Constitución y afirmó que el gobierno es de leyes y no de hombres, pese a ser el poder más frágil por no tener ni la bolsa ni la espada.
} \\
\end{longtable}

%==================================================================
% CAPÍTULO 2
%==================================================================
\chapter{Fallo Sojo c/ Poder Ejecutivo Nacional (1887) y competencia de la Corte}
\markboth{Fallo Sojo (1887) y competencia de la Corte}{}

\section{Introducción y contexto histórico argentino}

El fallo Sojo presenta una circunstancia similar a la de Marbury v. Madison, pero en Argentina. En cuanto al marco temporal: la Revolución de Mayo fue en 1810, la Independencia en 1816 y la primera Constitución en 1853. Para algunos la primera es la de 1860, porque Buenos Aires no estaba en la de 1853; esto es opinable.

La controversia del fallo Sojo se planteó en \textbf{1887}. Hacía poco que estaban en vigencia la división de poderes y la Constitución: veinticinco o veintisiete años es poco tiempo en la historia de un país. Como en Marbury v. Madison, hacía poco tiempo que regían la democracia y la Constitución, y había surgido un problema grande.

\section{Sojo: el caricaturista y la sátira política}

Sojo era un dibujante o caricaturista. En esa época no existía internet, y el medio eran los dibujos. Hacía \textbf{sátiras políticas} para criticar al poder a través de dibujos con ironía y sarcasmo, burlándose del poder.

El dibujo que motivó el caso representaba a una mujer con la cabeza sobre una mesada, con una prensa, una especie de bloque de hierro al que se le iban dando vueltas para aplastarle la cabeza. Alrededor de ella había seres grotescos, diabólicos y fantasmagóricos, con los bolsillos llenos de dinero, que se burlaban mientras le aplastaban la cabeza.

El significado del dibujo era el siguiente:
\begin{itemize}
    \item la \textbf{mujer} era la República Argentina;
    \item la \textbf{prensa}, lo que iba aplastando, era el sistema;
    \item los \textbf{seres siniestros, perversos y grotescos} representaban a los políticos que estaban robando, con los bolsillos llenos.
\end{itemize}

Es decir, la República languideciente y sufriente, y políticos alrededor, felices y grotescos, con carcajadas y bolsillos llenos, robándola.

\section{La detención arbitraria y el habeas corpus}

% TODO: contenido incompleto o ambiguo en las notas originales (el apunte atribuye la detención tanto al Poder Legislativo como al gobierno de Juárez Celman)
Ante el dibujo, el Poder Legislativo, indignado, se dio por aludido. El gobierno de Juárez Celman ordenó la detención de Sojo en prisión por el dibujo, al considerar que había dañado a una de las instituciones de la República. Quedó preso \textbf{sin juicio previo}, lo cual afecta el principio de inocencia.

\begin{definition}{Habeas corpus}
Figura jurídica muy antigua, que proviene de la época de la Gloriosa Revolución Inglesa, cuando se sostuvo que el monarca no tenía derecho a andar secuestrando gente. En latín significa \textit{``dónde está el cuerpo''}, \textit{``cuiden al cuerpo''}, \textit{``no dañen el cuerpo''}. Protege la libertad física de la persona.
\end{definition}

Frente a la detención, el abogado de Sojo presentó un habeas corpus.

\section{El error procesal: la instancia equivocada}

El abogado presentó el habeas corpus \textbf{directamente ante la Corte Suprema}. La pregunta es si hizo bien. En Argentina el sistema funciona principalmente por vía recursiva, por apelación, y no por acceso directo a la Corte, salvo en casos específicos, como los que involucran a estados o provincias, embajadores, etc.

El dilema del caso es por qué se lo presentó ante la Corte y no ante un juez ordinario, si todos los jueces controlan la Constitución. En esa época, como en la actual, los procedimientos y la división de poderes no eran del todo claros: era algo nuevo, y se estaba construyendo esa división tras las disputas entre unitarios y federales.

Para resolver el caso había que interpretar dos artículos concretos de la Constitución y la \textbf{Ley 48}, que ya existía en esa época y plantea cuáles son las cuestiones federales. La Ley 48 \textbf{sistematiza el control de constitucionalidad}, que está presente en dos artículos de la Constitución: el 116 y el 117.

\section{Artículo 116: la regla general}

\begin{quote}
\textit{``Corresponde a la Corte Suprema y a los tribunales inferiores de la Nación, el conocimiento y decisión de todas las causas que versen sobre puntos regidos por la Constitución y por las leyes de la Nación, con la reserva hecha en el inciso 12 del artículo 75; y por los tratados con las naciones extranjeras; de las causas concernientes a embajadores, ministros públicos y cónsules extranjeros; de las causas de almirantazgo y jurisdicción marítima; de los asuntos en que la Nación sea parte; de las causas que se susciten entre dos o más provincias; entre una provincia y los vecinos de otra; entre los vecinos de diferentes provincias; y entre una provincia o sus vecinos contra un Estado o ciudadano extranjero.''}
\end{quote}

\begin{note}
En el fallo Sojo este artículo se cita como artículo 100, ya que el número cambió con la Reforma Constitucional de 1994.
\end{note}

El artículo 116 establece el \textbf{género} referido al control de constitucionalidad. La regla general es que las causas se plantean por \textbf{vía recursiva o de apelación}: se pasa por la primera instancia, luego por la segunda y, si hay una causa, se llega a la Corte. Es decir, se deben agotar las instancias desde la más baja hasta la más alta. La última es la Corte, que interpreta, y no al revés.

Además, el artículo dice ``corresponde a la Corte Suprema y a los tribunales inferiores'': significa que corresponde a \textbf{ambos}. Lo que Marbury reclamaba era que el Poder Judicial los pusiera en el cargo, es decir, que completara el trámite, dado que el gobierno anterior sí había hecho el expediente de nombramiento.

En el caso de Sojo, por ser español y ciudadano extranjero, encajaba en el artículo 116. Podría haber ido por la vía recursiva después de la primera instancia, mediante el habeas corpus. Pero su abogado presentó la causa directamente ante la Corte.

\section{Artículo 117: la excepción, competencia originaria de la Corte}

\begin{quote}
\textit{``En estos casos [los del art. 116], la Corte Suprema ejercerá su jurisdicción por apelación según las reglas y excepciones que prescriba el Congreso; pero en todos los asuntos concernientes a embajadores, ministros y cónsules extranjeros, y en los que alguna provincia fuese parte, la ejercerá originaria y exclusivamente.''}
\end{quote}

\begin{note}
En el fallo Sojo este artículo se cita como artículo 101. El número cambió con la Reforma Constitucional de 1994.
\end{note}

El artículo 117 establece que en determinados casos sí se va \textbf{directamente a la Corte}. Esos casos están claramente determinados: embajadores, ministros y cónsules extranjeros, y los asuntos en que alguna provincia sea parte. % TODO: contenido incompleto o ambiguo en las notas originales (el docente enumera además casos entre provincias, entre una provincia y el Estado y contra ciudadanos extranjeros, que en el texto citado corresponden al art. 116)
Sojo, aunque era español y extranjero, no era cónsul ni embajador ni nada similar: era un simple dibujante.

\begin{important}
La competencia originaria de la Corte es \textbf{excepcional, no es la regla}. La regla es la competencia ordinaria, y excepcionalmente la extraordinaria, que tiene ciertos requisitos.
\end{important}

\begin{definition}{Numerus clausus}
Expresión latina que significa \textit{``número cerrado''}. Los casos de competencia originaria de la Corte del artículo 117 están taxativamente determinados: si un caso no está en el artículo 117, no es ampliable ni hay forma de decir que otro entra. Son estos casos y no otros; no se admite interpretación extensiva ni la incorporación de otros supuestos por analogía.
\end{definition}

\begin{table}[H]
\centering
\begin{tabularx}{\textwidth}{
    >{\raggedright\arraybackslash}p{0.20\textwidth}
    >{\raggedright\arraybackslash}p{0.17\textwidth}
    >{\raggedright\arraybackslash}X
    >{\raggedright\arraybackslash}p{0.13\textwidth}
}
\toprule
\textbf{Tipo de competencia} & \textbf{Vía de acceso} & \textbf{Casos} & \textbf{Norma} \\
\midrule
\textbf{Ordinaria (regla)} & Vía de apelación o recursiva. & Todas las causas de naturaleza federal. & Art. 116 CN \\
\addlinespace
\textbf{Originaria y exclusiva (excepción)} & Acceso directo a la Corte. & Embajadores, cónsules, ministros extranjeros y causas donde alguna provincia sea parte. & Art. 117 CN \\
\bottomrule
\end{tabularx}
\caption{Competencia ordinaria y competencia originaria de la Corte.}
\end{table}

\subsection{Ejemplo práctico: el accidente del embajador de Chile}

\begin{example}{Accidente de tránsito con un embajador}
Se escucha una explosión fuerte, como de un choque gigante: un auto chocó sobre el Libertador y Pueyrredón. Es un accidente de tránsito. ¿Qué tipo de juez y de qué instancia interviene? Un juez civil de primera instancia.

Ahora bien, quien iba dentro del auto era el Embajador de Chile. ¿Sigue siendo competente el juez civil de primera instancia? No. A raíz del artículo 117, por la razón de la persona, es decir, por la calidad de cónsul o embajador, el caso es de \textbf{competencia originaria de la Corte}.
\end{example}

Sojo no era embajador ni cónsul: era un dibujante. Por lo tanto, la competencia originaria de la Corte, que es excepcional, no era aplicable.

\section{La decisión de la Corte en el fallo Sojo}

La Corte interpretó que no se trataba de una cuestión de inconstitucionalidad sino de \textbf{jurisdicción y competencia}. El caso no tendría que haber llegado a la Corte, sino que debió tramitar ante el juez de primera instancia; si allí se violaban determinadas cuestiones procedimentales, podría tener la posibilidad de llegar después a la Corte. Por ello la Corte no podía tomar el caso: Sojo se había equivocado de instancia.

La Corte remitió la causa al juez de primera instancia, quien aplicó el habeas corpus y dejó libre a Sojo. La trascendencia de esta solución es clara: de haberse limitado a decir que Sojo se equivocó de instancia y que estaba bien juzgado, Sojo habría permanecido detenido.

\begin{table}[H]
\centering
\begin{tabularx}{\textwidth}{
    >{\raggedright\arraybackslash}p{0.20\textwidth}
    >{\raggedright\arraybackslash}X
    >{\raggedright\arraybackslash}X
}
\toprule
& \textbf{Marbury v. Madison} & \textbf{Sojo c/ Poder Ejecutivo Nacional} \\
\midrule
\textbf{Cuestión de fondo} & Declaración de inconstitucionalidad. & No hay cuestión de inconstitucionalidad, sino aplicación correcta de la jurisdicción y competencia. \\
\addlinespace
\textbf{Vía elegida} & Directamente ante la Corte. & Directamente ante la Corte. \\
\addlinespace
\textbf{Resultado} & La Corte no nombra a los jueces, pero se consolida como único intérprete de la Constitución. & La Corte remite el caso a primera instancia, donde se aplica el habeas corpus y Sojo queda libre. \\
\bottomrule
\end{tabularx}
\caption{Comparación entre los fallos Marbury y Sojo.}
\end{table}

\section{Cuadro Cornell: Sojo y los artículos 116 y 117}

\begin{longtable}{@{}>{\RaggedRight\arraybackslash}p{0.27\textwidth}>{\RaggedRight\arraybackslash}p{0.64\textwidth}@{}}
\toprule
\textbf{Preguntas / palabras clave} & \textbf{Notas / respuestas} \\
\midrule
\endfirsthead
\toprule
\textbf{Preguntas / palabras clave} & \textbf{Notas / respuestas} \\
\midrule
\endhead
\bottomrule
\endfoot
\bottomrule
\endlastfoot

\textbf{¿Quién era Sojo y por qué fue detenido?} &
Era un dibujante o caricaturista que hacía sátira política. Fue detenido en 1887 sin juicio previo por un dibujo que representaba a la República Argentina, simbolizada por una mujer, aplastada por el sistema mientras políticos siniestros con los bolsillos llenos se burlaban de ella. Se consideró que había dañado a una de las instituciones de la República.
\cornellsep

\textbf{¿Qué es el habeas corpus y cuál es su origen?} &
Es una figura muy antigua, que remite a la Gloriosa Revolución Inglesa, cuando el monarca no tenía derecho a secuestrar gente. Significa ``dónde está el cuerpo'' y protege la libertad física de la persona.
\cornellsep

\textbf{¿Cuál fue el error procesal en el caso Sojo?} &
Su abogado presentó el habeas corpus directamente ante la Corte Suprema, en lugar de hacerlo ante un juez de primera instancia. Al ser español y extranjero, Sojo encajaba en el artículo 116 y podría haber ido por la vía recursiva después de la primera instancia; pero la regla es la vía de apelación y no el acceso directo.
\cornellsep

\textbf{¿Qué rol cumple la Ley 48?} &
Plantea cuáles son las cuestiones federales y sistematiza el control de constitucionalidad, que está presente en los artículos 116 y 117 de la Constitución.
\cornellsep

\textbf{¿Qué establece el artículo 116 de la Constitución Nacional?} &
Es la regla general del control de constitucionalidad. Las causas que versan sobre puntos regidos por la Constitución, las leyes de la Nación y los tratados corresponden a la Corte Suprema y a los tribunales inferiores, y se plantean por vía de apelación, agotando las instancias desde la más baja hasta la más alta. La Corte interpreta en último lugar.
\cornellsep

\textbf{¿Qué establece el artículo 117 y en qué se diferencia del 116?} &
Es la excepción: la Corte ejerce competencia originaria y exclusiva, es decir, se accede directamente a ella, en los asuntos concernientes a embajadores, ministros y cónsules extranjeros y en los que alguna provincia sea parte. El artículo 116 establece la regla (vía de apelación); el 117 establece la excepción.
\cornellsep

\textbf{¿Qué significa numerus clausus?} &
Significa ``número cerrado'': los casos de competencia originaria del artículo 117 son taxativos y no ampliables. Solo se va directamente a la Corte en los casos allí enumerados, y no hay forma de agregar otros.
\cornellsep

\textbf{¿Por qué el accidente del embajador de Chile es competencia de la Corte y un accidente común no?} &
Un accidente de tránsito común corresponde a un juez civil de primera instancia. Si el conductor es el Embajador de Chile, por el artículo 117 y por la calidad de la persona, el caso pasa a ser de competencia originaria de la Corte.
\cornellsep

\textbf{¿Cómo resolvió la Corte el caso Sojo?} &
Consideró que no era una cuestión de inconstitucionalidad sino de jurisdicción y competencia: Sojo se había equivocado de instancia. Remitió la causa al juez de primera instancia, quien aplicó el habeas corpus y lo dejó libre. En Marbury, en cambio, la cuestión era la declaración de inconstitucionalidad. \\
\midrule

\multicolumn{2}{@{}p{\dimexpr0.91\textwidth+2\tabcolsep\relax}@{}}{
\textbf{Resumen:} Sojo, dibujante detenido sin juicio previo por una caricatura política, presentó un habeas corpus directamente ante la Corte Suprema. El artículo 116 establece como regla el acceso por vía de apelación, agotando instancias, y el artículo 117 establece como excepción, en un numerus clausus, la competencia originaria de la Corte para embajadores, cónsules, ministros extranjeros y causas donde una provincia sea parte. Como Sojo no encuadraba en esa excepción, la Corte declaró que se había equivocado de instancia y remitió el caso a primera instancia, donde fue liberado.
} \\
\end{longtable}

%==================================================================
% CAPÍTULO 3
%==================================================================
\chapter{Control difuso y Recurso Extraordinario Federal}
\markboth{Control difuso y Recurso Extraordinario Federal}{}

\section{El sistema difuso de control de constitucionalidad}

\begin{definition}{Control difuso de constitucionalidad}
Sistema, adoptado por Argentina y por los Estados Unidos, en el que todos los jueces, y no solo la Corte Suprema, tienen el poder de declarar inconstitucional una norma en un caso concreto. Se diferencia del sistema concentrado, propio del modelo europeo o francés, en el que solo un tribunal constitucional especializado puede hacerlo.
\end{definition}

Por eso en el fallo Sojo se plantea que todos los jueces controlan la Constitución, y por eso la Corte Suprema no es el único órgano que interviene en las cuestiones constitucionales, aunque actúa como árbitro último.

\section{El Recurso Extraordinario Federal (REF)}

\subsection{El mecanismo para llegar a la Corte Suprema}

\begin{definition}{Recurso Extraordinario Federal (REF)}
Instrumento procesal específico para acceder a la Corte Suprema de Justicia de la Nación. No se utiliza otro mecanismo para llegar a la Corte.
\end{definition}

Para que la Corte tenga la posibilidad de tomar una causa, esta debe tener \textbf{naturaleza federal}. Si una cuestión no tiene naturaleza federal, no existe alternativa ni posibilidad de llegar a la Corte. Esto no significa una negación de justicia: la Corte Suprema está para el tratamiento de cuestiones federales y no de cuestiones comunes. Para estas últimas están los demás jueces, y por eso el control es difuso, a diferencia del sistema francés.

La Corte se aboca a cuestiones específicas: de \textbf{gravedad institucional}, de interpretación constitucional profunda, de tratados o de relevancia constitucional.

\subsection{Planteo versus reserva de la cuestión federal}

En la práctica forense está muy difundida la expresión ``reservo la cuestión federal'' o ``reservo caso federal'', que aparece en los escritos e incluso es utilizada por docentes. El docente sostiene que esa práctica es incorrecta y reconoce que su posición se aparta de la habitual: las cuestiones federales \textbf{se plantean y no se reservan}. En el mundo del derecho las cuestiones no se reservan, se ejercen.

\begin{important}
\textbf{Los derechos se ejercen, no se reservan.} Según el docente, la cuestión federal debe plantearse expresamente en el proceso, no meramente reservarse. La expresión ``reservo la cuestión federal'' es técnicamente incorrecta y podría llevar al rechazo del escrito.
\end{important}

\section{Cuadro Cornell: control difuso y REF}

\begin{longtable}{@{}>{\RaggedRight\arraybackslash}p{0.27\textwidth}>{\RaggedRight\arraybackslash}p{0.64\textwidth}@{}}
\toprule
\textbf{Preguntas / palabras clave} & \textbf{Notas / respuestas} \\
\midrule
\endfirsthead
\toprule
\textbf{Preguntas / palabras clave} & \textbf{Notas / respuestas} \\
\midrule
\endhead
\bottomrule
\endfoot
\bottomrule
\endlastfoot

\textbf{¿Qué es el control difuso de constitucionalidad?} &
Es el sistema, adoptado por Argentina y por los Estados Unidos, en el que todos los jueces pueden declarar inconstitucional una norma en un caso concreto. Se diferencia del sistema concentrado (modelo europeo o francés), en el que solo un tribunal constitucional especializado puede hacerlo.
\cornellsep

\textbf{¿Qué es el Recurso Extraordinario Federal (REF)?} &
Es el instrumento procesal específico para acceder a la Corte Suprema de Justicia de la Nación; no se utiliza otro mecanismo. Para que la Corte pueda tomar la causa, esta debe tener naturaleza federal.
\cornellsep

\textbf{¿Por qué se exige naturaleza federal y a qué cuestiones se aboca la Corte?} &
Porque la Corte Suprema está para el tratamiento de cuestiones federales y no de cuestiones comunes, que corresponden a los demás jueces; esto no implica negación de justicia. La Corte se aboca a cuestiones de gravedad institucional, interpretación constitucional profunda, tratados o relevancia constitucional.
\cornellsep

\textbf{¿Se ``reserva'' o se ``plantea'' la cuestión federal?} &
Según el docente, la postura técnicamente correcta es que las cuestiones federales se plantean y no se reservan. La expresión ``reservo la cuestión federal'', aunque muy difundida en la práctica forense, es cuestionada porque los derechos se ejercen, no se reservan. Según el apunte, plantear la cuestión federal es condición de admisibilidad del REF. \\
\midrule

\multicolumn{2}{@{}p{\dimexpr0.91\textwidth+2\tabcolsep\relax}@{}}{
\textbf{Resumen:} El control de constitucionalidad en el sistema difuso, adoptado por Argentina, es ejercido por todos los jueces, pero es la Corte Suprema quien actúa como árbitro último. Se accede a ella mediante el Recurso Extraordinario Federal, cuyo requisito ineludible es la naturaleza federal de la cuestión y su planteo expreso, no su reserva, en el proceso.
} \\
\end{longtable}

\end{document}
